% =====================================================================
% Modeling Factory -- MCM/ICM paper template (Step 9 starting point)
% =====================================================================
% Copy this file to <project>/<base>_paper.tex and fill it in; never
% \input it from the factory.  It compiles as-is with compile_paper.sh.
%
% Class: mcmthesis v6.3.3 (LPPL 1.3c), vendored unmodified in
% __FACTORY__/latex_templates/ next to this file.  compile_paper.sh puts
% that directory on TEXINPUTS, so the project needs no copy of the class
% and `\usepackage{abstract_placeholder}` resolves the same way.
%
% COMAP rules this layout follows:
%   * Page 1 is the one-page Summary Sheet (problem, year, control number,
%     title, summary, keywords).  The summary must fit on that page.
%   * At most 25 pages from the Summary Sheet through the appendices
%     (memo/letter, references and code appendices count).
%   * The Report on Use of AI Tools follows \AImatter and does not count
%     toward the 25 pages; it is required when AI tools were used, and the
%     tools are also cited inline and listed in the references (COMAP AI
%     policy v102025).
%   * No names of team members, advisor or school on any page; the header
%     carries only the Team Control Number.
%   * English, readable font of at least 12 pt (the class is 12 pt).
% =====================================================================
\documentclass{mcmthesis}
\mcmsetup{
  tcn = {0000000},   % Team Control Number: placeholder, set at submission
  problem = {A},     % Problem Chosen: A-F, from problem/problem_brief.md
  sheet = true, titleinsheet = true, keywordsinsheet = true,
  titlepage = false, abstract = false}
% The class prints \the\year (the compile date) on the Summary Sheet; pin
% the contest year so a later recompile of unchanged sources cannot change it.
\renewcommand{\headset}{2027\\MCM/ICM\\Summary Sheet}
\setlength{\headheight}{14pt}  % room for the "Team # ... Page X of Y" header
% COMAP's AI policy names the appendix "Report on Use of AI Tools".
\renewcommand{\reportname}{Report on Use of AI Tools}

\usepackage{newtxtext}   % Times-like text face (COMAP asks for 12 pt Times)
\usepackage{newtxmath}
\usepackage{float}
\usepackage{caption}
\usepackage{subcaption}
\usepackage{enumitem}
\graphicspath{{figures/}}
\usepackage{abstract_placeholder}

\title{<Informative title: object + main method, about 10--15 words>}
\author{}
\date{}

\begin{document}

% ======================= Summary Sheet (page 1) =======================
% Step 9 leaves the placeholder line below untouched; Step 14 replaces it
% with the English summary (one page together with the keywords).
\begin{abstract}
\AbstractPlaceholder % ABSTRACT_PLACEHOLDER compatibility marker

\begin{keywords}
<Keyword 1>; <Keyword 2>; <Keyword 3>; <Keyword 4>
\end{keywords}
\end{abstract}
\maketitle

\tableofcontents
\newpage

% ============================ Introduction ============================
\section{Introduction}
\subsection{Problem Background}
% Restate the context in your own words; do not copy the problem text.
\subsection{Restatement of the Problem}
% One item per task of the problem statement, in the problem's order.
\subsection{Our Work}
% One paragraph plus the workflow figure (an explain_model figure from
% visualization_log.md) mapping each task to the model that answers it.

% ===================== Assumptions and Justifications =================
\section{Assumptions and Justifications}
% From assumption_ledger.md; each assumption with its justification.
\begin{enumerate}[label=\textbf{Assumption \arabic*:}, leftmargin=*]
  \item <Statement>.
  \textbf{Justification:} <one sentence>.
\end{enumerate}

% ============================== Notations =============================
\section{Notations}
% Must agree with symbol_table.md.
\begin{table}[H]
  \centering
  \caption{Notations}
  \begin{tabular}{cll}
    \toprule
    Symbol & Description & Unit \\
    \midrule
    $x_i$ & <description> & <unit> \\
    \bottomrule
  \end{tabular}
\end{table}

% ======================= Data Preprocessing (opt.) ====================
% \section{Data Preprocessing}

% ===================== Models: one section per task ===================
\section{Model I: <Task 1 model name>}
\subsection{Model Formulation}
\subsection{Solution}
\subsection{Results}
% Lead with the final answer and its main figure or table, then the
% derivation and the checks.

\section{Model II: <Task 2 model name>}
\subsection{Model Formulation}
\subsection{Solution}
\subsection{Results}

% ====================== Sensitivity Analysis ==========================
\section{Sensitivity Analysis}
% From sensitivity_report.md.

% ================= Model Evaluation and Further Discussion ============
\section{Model Evaluation and Further Discussion}
\subsection{Strengths}
\subsection{Weaknesses}
\subsection{Further Discussion}
% From evaluation.md.

% ============================== Conclusion ============================
\section{Conclusion}

% ==================== Memo / Letter (only if required) =================
% Keep only when the problem asks for a memo, letter, policy brief or
% similar non-technical document (common in Problems C--F).  It counts
% toward the 25 pages and is usually limited to one or two pages.
% \memoto{<Recipient>}
% \memofrom{Team \#\,0000000}
% \memosubject{<Subject>}
% \memodate{<Date>}
% \begin{memo}[Memorandum]
% <Body>
% \end{memo}

% ============================== References ============================
\bibliographystyle{plain}
\bibliography{references}

% ============================== Appendices ============================
\begin{appendices}
\section{Core Code}
% Only the code a reader needs to reproduce the main results; trim it
% first when the page count is over 25.
% \lstinputlisting[language=Python]{models/<stream>/03_solve.py}
\end{appendices}

% ========= Report on Use of AI Tools (not counted in 25 pages) ========
% \AImatter closes the counted pages (the header "Page X of Y" ends here)
% and switches to the uncounted AI report.  List every AI tool that
% produced content in this paper: model name and version, what it was
% used for, and the queries with their outputs (summarised when long).
% The human team confirms this report before submission.
\AImatter
\begin{ReportAiUse}{9}
  \bibitem{AI1} <Model name and version, provider>\\
  Query1: <query>\\
  Output: <output>
\end{ReportAiUse}

\end{document}
